\newpage
\section{Grundlagen}

\subsection{Large Language Models}

 \acp{LLM} sind maschinelle Lernmodelle, die auf sehr großen Textmengen trainiert werden, um natürliche Sprache zu verarbeiten und zu erzeugen. Technisch handelt es sich um neuronale Netze mit typischerweise mehreren hundert Milliarden Parametern \parencite[S.~4]{zhao_survey_2026}. Dadurch sind sie in der Lage, komplexe Muster, sprachliche Semantik und Nuancen zu erlernen, die dann für viele verschiedene sprachbezogene Aufgaben wie Übersetzung, Textzusammenfassungen etc. verwendet werden können \parencite[S.~2]{celik_review_2025}. Den Durchbruch verdanken \acp{LLM} durch die von \textcite{vaswani_attention_2017} eingeführte Transformer-Architektur \parencite[S.~2]{celik_review_2025}.

\subsubsection{Funktionsprinzip und Kontext als Qualitätsfaktor}

Der zentrale Bestandteil der Transformer-Architektur ist der Self-Attention-Mechanismus. Er setzt die Elemente einer Eingabe direkt miteinander in Beziehung, unabhängig davon, wie weit sie voneinander entfernt stehen \parencite[S.~6]{vaswani_attention_2017}. Frühere rekurrente Modelle mussten die Eingabe dafür Schritt für Schritt durchlaufen, was eine Parallelisierung innerhalb eines Trainingsbeispiels ausschloss. Der Transformer verzichtet auf diese Rekurrenz, verarbeitet die Eingabe weitgehend parallel und verkürzt damit die Trainingszeit erheblich \parencite[S.~1]{vaswani_attention_2017}. 

Vor der Verarbeitung wird der Text in Tokens zerlegt, also in kleinere Einheiten wie Wörter oder Teilwörter, damit sie von \acp{LLM} einfacher und besser verarbeitet werden können \parencite[S.~18]{hou_large_2024}. Ein verbreitetes Verfahren zur Zerlegung in solche Teilwörter ist das Byte-Pair-Encoding (BPE), das ursprünglich für die maschinelle Übersetzung entwickelt wurde \parencite[S.~2]{sennrich_neural_2016} und seitdem auch bei \acp{LLM} zur Tokenisierung eingesetzt wird \parencite[S.~19]{zhao_survey_2026}. Die Texterzeugung erfolgt autoregressiv. Bei jedem Schritt verwendet das Modell die bereits erzeugten Tokens als zusätzliche Eingabe und berechnet eine Wahrscheinlichkeitsverteilung für das nächste Token \parencite[S.~2,~5]{vaswani_attention_2017}. Wie stark die Auswahl vom jeweils wahrscheinlichsten Token abweichen darf, steuert unter anderem der Temperatur-Parameter. Er bestimmt den Grad an Zufälligkeit im erzeugten Text \parencite[S.~4]{plein_automatic_2023}.

\textcite[S.~1,~4]{brown_language_2020} zeigten, dass große Sprachmodelle Aufgaben allein anhand einer Anweisung oder einzelner Beispiele im Prompt lösen können, ohne dass die Modellgewichte angepasst werden, und bezeichnen dies als \enquote{in-context-learning}. Ob das Modell dabei im eigentlichen Sinne lernt oder lediglich im Training gesehene Muster wiedererkennt, lassen die Autoren offen \parencite[S.~34]{brown_language_2020}. Für diese Arbeit genügt die Beobachtung, dass im Prompt bereitgestellte Informationen die Ausgabe unmittelbar bestimmen.

Daraus folgt ein zentraler Punkt für diese Arbeit. Ein LLM hat kein konkretes Wissen über eine individuell entwickelte Benutzeroberfläche. Informationen über die tatsächlich vorhandenen \ac{UI}-Elemente, deren Selektoren und Zustände stehen weder in den Trainingsdaten noch sind sie zur Laufzeit zugänglich, solange sie nicht explizit im Prompt bereitgestellt werden. Für die Testgenerierung folgt daraus, dass Qualität und Umfang des bereitgestellten Kontexts den generierten Testcode maßgeblich mitbestimmen dürften. Genau dieser Zusammenhang wird in dieser Arbeit untersucht

\subsubsection{Einsatz im Software Engineering}

\acp{LLM} werden zunehmend im Software Engineering eingesetzt, insbesondere zur Generierung von Quellcode. Sichtbar wird das an Werkzeugen wie GitHub Copilot, die auf Grundlage des umgebenden Quellcode-Kontexts Vorschläge zur Codevervollständigung erzeugen \parencite[S.~30]{hou_large_2024}. Der systematische Literaturüberblick von \textcite{hou_large_2024} beschreibt 85 konkrete Aufgaben aus sechs Kernaktivitäten des Software Engineering, von der Anforderungsanalyse bis zur Wartung, und belegt damit, dass das Feld deutlich über die Codevervollständigung hinausreicht \parencite[S.~4]{hou_large_2024}.

Ein Anwendungsfeld ist die automatisierte Testgenerierung, in der aus natürlichsprachlichen oder strukturierten Beschreibungen Tests auf unterschiedlichen Teststufen erzeugt werden \parencite{celik_review_2025,ferreira_acceptance_2025,alian_feature-driven_2024}. Diesen Möglichkeiten stehen bekannte Grenzen gegenüber. Da \acp{LLM} ihre Ausgabe auf Basis von Wahrscheinlichkeiten erzeugen, können sie plausibel wirkende, faktisch aber nicht existierende Inhalte produzieren. \textcite{celik_review_2025} fassen zusammen, dass sich über die von ihnen betrachteten Studien zum Prompt Engineering hinweg wiederkehrende Probleme zeigen. Dazu gehören halluzinierter Code, Kompilierfehler und eine eingeschränkte Übertragbarkeit auf reale Umgebungen, besonders dort, wo der Projektkontext dünn oder unvollständig ist \parencite[S.~10]{celik_review_2025}.

Grundsätzlich lassen sich Sprachmodelle auf zwei Wegen an eine spezifische Aufgabe anpassen. Beim Fine-Tuning werden die Gewichte eines vortrainierten Modells durch Training auf einen aufgabenspezifischen, gelabelten Datensatz angepasst, wofür typischerweise mehrere Tausend bis mehrere Hunderttausend Beispiele nötig sind \parencite[S.~6]{brown_language_2020}. Beim Prompting wird das Verhalten allein über die Gestaltung der Eingabe gesteuert. Es beruht auf dem in Abschnitt~2.1.1 beschriebenen In-Context-Learning, das in dieser Arbeit verfolgt wird. Nachteile des Fine-Tuning sind der Bedarf an einem großen Datensatz für jede Aufgabe und eine schlechte Übertragbarkeit außerhalb der Trainingsdaten \parencite[S.~6]{brown_language_2020}. Ein prompt-basierter Workflow lässt sich jedoch ohne erneutes Training auf andere Anwendungen und Use Cases übertragen, während ein feinjustiertes Modell auf seinen Trainingskontext beschränkt bliebe.

\subsection{E2E-Testautomatisierung mit Playwright}

\subsubsection{Teststufen und Einordnung von E2E-Tests}

Softwaretests lassen sich nach ihrem Umfang und ihrer Nähe zum Nutzerverhalten in verschiedene Stufen einteilen. Eine verbreitete Einteilung von Softwaretests ist die Testpyramide, die von \textcite{poenisch_testpyramide_2022} beschrieben wird. Sie unterscheidet drei Ebenen, wie in Abbildung~\ref{fig:testpyramide} dargestellt. Ganz unten stehen Unit-Tests, die einzelne Funktionen oder Komponenten prüfen. Sie laufen schnell und werden in großer Anzahl geschrieben. Darüber liegen die Integrationstests, die das Zusammenspiel mehrerer Komponenten betrachten. An der Spitze stehen Systemtests, die das Verhalten des gesamten Systems betrachten. Zu ihnen gehören \ac{UI}-Tests und \ac{E2E}-Tests. Da dabei ausschließlich das äußere Verhalten der Anwendung geprüft wird, handelt es sich um Black-Box-Tests, welche die Anwendung als Ganzes aus Sicht der Nutzer betrachten.

\begin{figure}[H]
    \centering
    \caption[Testpyramide nach Mike Cohn]{Testpyramide nach Mike Cohn mit den drei Teststufen Unit-, Integrations- und Systemtests. Die Breite einer Ebene steht für die empfohlene Anzahl an Tests: viele schnelle Unit-Tests an der Basis, wenige langsame Systemtests an der Spitze. Die in dieser Arbeit generierten \ac{E2E}-Tests gehören zur obersten Ebene. (Quelle: \cite{poenisch_testpyramide_2022})}
    \label{fig:testpyramide}
    \includegraphics[width=0.7\textwidth]{testpyramide}
\end{figure}

Ein \ac{E2E}-Test steuert dazu einen Browser und führt dieselben Schritte aus wie ein Mensch. Dazu gehören das Klicken auf Schaltflächen, das Ausfüllen von Feldern und das Navigieren durch die Oberfläche. Geprüft wird, ob am Ende das erwartete Ergebnis eintritt. Dadurch geben \ac{E2E}-Tests eine gute Auskunft darüber, ob eine fachliche Anforderung tatsächlich umgesetzt ist. Allerdings sind sie langsam, aufwendig zu erstellen und vergleichsweise anfällig und können schon bei kleinen Änderungen an der \ac{UI} fehlschlagen \parencite{fowler_test_2012}. Deshalb wird empfohlen, ihre Anzahl, entsprechend der Form der Pyramide, gering zu halten \parencite{poenisch_testpyramide_2022}.

Für WebGIS-Anwendungen sind \ac{E2E}-Tests besonders relevant. Viele Probleme zeigen sich erst im Zusammenspiel der Komponenten und Interaktionen. Dazu gehört das asynchrone Laden von Geodaten, das Rendern der Karte oder die Auswirkungen der Layersteuerung auf die Kartenanzeige. Solche Abläufe lassen sich mit Unit-Tests kaum abbilden, weil sie den vollständigen, im Browser laufenden Zustand voraussetzen.

\subsubsection{Playwright: Selektoren, Assertions, asynchrone UIs}

Playwright ist ein Open-Source-\ac{E2E}-Test-Framework von Microsoft zur Automatisierung von Browsern. Es steuert die Browser-Engines Chromium, Firefox und WebKit über eine einheitliche Programmierschnittstelle und führt Tests in einem echten Browser aus \parencite[Playwright]{microsoft_fast_nodate}. Die Tests können in JavaScript, TypeScript, Python, Java und .NET geschrieben werden \parencite[Supported languages]{microsoft_supported_nodate}.

Um mit der Benutzeroberfläche zu interagieren, muss ein Test die jeweiligen Elemente finden. Playwright nutzt dafür sogenannte Lokatoren. Es wird empfohlen, die Elemente über nutzerseitige Eigenschaften, wie über ihre \ac{ARIA}-Rolle oder ihren sichtbaren Text (\code|getByText|, \code|getByRole|), anzusprechen. Solche Selektoren sind robuster gegenüber Änderungen am Markup als \ac{CSS}- oder XPath-Ausdrücke, die an die konkrete \ac{DOM}-Struktur gebunden sind. Playwright kann diese Informationen auch gebündelt aus einem Accessibility-Snapshot auslesen. Er gibt Rollen, Attribute, Werte und Textinhalt der Elemente einer Seite als Baumstruktur aus, damit man eine Übersicht erhält, welche Elemente alles vorhanden sind \parencite[Snapshot testing]{microsoft_snapshot_nodate}. Eine Alternative sind Test-IDs. Über das Attribut \code|data-testid| lässt sich ein Element gezielt für Tests markieren und mit \code|getByTestId| ansprechen, unabhängig von Darstellung und Position. Welches Attribut als Test-ID gilt, ist konfigurierbar \parencite[Locators]{microsoft_locators_nodate}.

Ein wesentliches Merkmal von Playwright ist das automatische Warten. Vor jeder Aktion prüft das Framework, ob das Zielelement sichtbar, stabil und bedienbar ist, ansonsten wird bis zu einem Timeout gewartet \parencite[Auto-waiting]{microsoft_auto-waiting_nodate}. Dadurch müssen nicht für jede Interaktion manuelle Wartezeiten eingebaut werden, was eine häufige Ursache für instabile Tests beseitigt. Dasselbe Prinzip gilt für Assertions. Mit \code|expect| formulierte Zusicherungen wie \code|expect(locator).toBeVisible()| oder \code|toHaveText()| werden so lange wiederholt, bis die Bedingung erfüllt ist oder die Zeit abläuft \parencite[Assertions]{microsoft_assertions_nodate}. Für Werte, die nicht von einem Lokator stammen, sondern z.~B. von einer eigenen Funktion zurückgegeben werden, lässt sich mit \code|expect.poll| auch wiederholt prüfen \parencite[expect.poll]{microsoft_assertions_nodate}. Eine Assertion prüft damit nicht einen Momentzustand, sondern einen Zustand, der sich auch verzögert einstellen darf.

Für Fälle, die das automatische Warten nicht abdeckt, stellt Playwright extra Wartemethoden bereit, etwa auf eine bestimmte Netzwerkantwort (\code|waitForResponse|) oder einen Ladezustand der Seite (\code|waitForLoadState|) \parencite[Page]{microsoft_page_nodate}. Das ist gerade bei WebGIS-Anwendungen wichtig, in denen Kartenkacheln und Geodaten asynchron über das Netz geladen werden und erst danach sichtbar auftauchen. Ein Test, der zu früh prüft, würde fehlschlagen, obwohl die Anwendung korrekt arbeitet. Gerade beim asynchronen Laden von Kartenkacheln und Geodaten greift damit beides ineinander: das automatische Warten für \ac{DOM}-Elemente und gezielte Wartepunkte auf Netzwerkebene.

Neben der Interaktion mit dem \ac{DOM} kann Playwright mit \code|page.evaluate()| auch beliebigen JavaScript-Code im Browser-Kontext ausführen und dabei auf globale Objekte der Seite zugreifen \parencite[Evaluating JavaScript]{microsoft_evaluating_nodate}. Das ermöglicht Zugriffe, die über die reine \ac{DOM}-Interaktion hinaus gehen.

\subsection{Open Pioneer Trails}

\subsubsection{Framework-Überblick und Komponentenmodell}

Open Pioneer Trails ist ein Open-Source-Framework zur Entwicklung webbasierter GeoIT-Anwendungen. Es ist als gemeinsames Projekt der con terra GmbH und 52°North Spatial Information Research GmbH im Rahmen der Open-Pioneer-Initiative entstanden und steht unter der Apache-2.0-Lizenz frei auf GitHub zur Verfügung. Trails ist ein Entwicklungsframework und kein fertiges Produkt. Anwendungen werden also nicht konfiguriert, sondern programmiert, wobei das Framework wiederverwendbare Komponenten und Vorlagen bereitstellt \parencite{52north_gmbh_conquering_2026}.

Technisch besteht Trails aus modernen Webtechnologien. Die Anwendungen werden in TypeScript mit React geschrieben, das Styling basiert auf Chakra UI, als Kartenkomponente wird OpenLayers verwendet, die Paketverwaltung läuft über pnpm (Performant Node Package Manager) und als Build-System wird Vite verwendet \parencite[trails-starter]{open_pioneer_open-pioneertrails-starter_2023}.

Eine Open Pioneer Trails-Anwendung ist auf Paketen aufgebaut. Ein Paket bündelt React-Komponenten, Services, Übersetzungen und Styling. Die Anwendung selbst ist ebenfalls ein Paket und wird zu einer Web Component gebaut, dadurch kann sie einfach in bestehende Web-Anwendungen integriert werden. Aktuell werden zwei Gruppen wiederverwendbarer Pakete bereitgestellt. Die Core Packages enthalten kartenlose Basisfunktionen wie Authentifizierung, Notifier, Internationalisierungsfunktionen und einen \ac{HTTP}-Service \parencite[trails-core-packages]{open_pioneer_open-pioneertrails-core-packages_2023}. Die OpenLayers Base Packages bauen darauf auf und liefern die Kartenkomponenten, darunter Layerübersicht, Kartennavigations-Werkzeuge, Legende, Selektion und Ergebnisliste, Suche, Messwerkzeug, Koordinatenanzeige und Kartendruck \parencite[trails-openlayers-base-packages]{open_pioneer_open-pioneertrails-openlayers-base-packages_2023}. Abbildung~\ref{fig:opt_oekosystem} gibt einen Überblick über das Paketökosystem.

\begin{figure}[H]
    \centering
    \caption[Open Pioneer Trails Paketökosystem]{Paketökosystem von Open Pioneer Trails. Die Core Packages stellen kartenlose Basisfunktionen bereit, die OpenLayers Base Packages bauen darauf auf und liefern die Kartenkomponenten. Projektspezifische Anwendungen nutzen beide Ebenen in unterschiedlichem Umfang.}
    \label{fig:opt_oekosystem}
    \includegraphics[width=0.4\textwidth]{opt_oekosystem_pakete}
\end{figure}

Die Anbindung der Karte erfolgt nicht direkt über OpenLayers, sondern über eine eigene Schnittstelle. Diese erweitert das Datenmodell, vereinfacht den Umgang mit dem Kartenzustand und stellt ihn als TypeScript-Schnittstelle (MapModel) bereit. Der direkte Zugriff auf die OpenLayers-\ac{API} bleibt bei Bedarf möglich \parencite{52north_gmbh_open_2025}.

Trails richtet sich explizit an Entwickler und erfordert Programmierkenntnisse. Dafür lässt es Layout und Funktionalität frei gestalten und bringt typische \ac{GIS}-Widgets als fertige Komponenten mit.

\subsubsection{data-testid-Konventionen und Testbarkeit}

Trails reicht \code|data-testid| über CommonComponentProps \parencite[react-utils]{open_pioneer_open-pioneerreact-utils_nodate} bzw. den Chakra-Prop-Filter \parencite{chakra_systems_chakra_2026} ans \ac{DOM} weiter.

Neben den Test-IDs sind die Komponenten auf Barrierefreiheit ausgelegt und folgen \ac{ARIA}-Konventionen \parencite{52north_gmbh_open_2025}. Elemente lassen sich dadurch beispielsweise über ihre Rolle oder ihren sichtbaren Textinhalt finden. Playwright kann anhand der Rolle und Bezeichnung (\code|getByRole|/\code|getByText|) auf diese Elemente zugreifen. Damit bietet Trails zwei robuste Wege, ein Element zu adressieren.

\subsection{WebGIS-Anwendungen und Testkomplexität}

WebGIS-Anwendungen stellen die Testautomatisierung vor Herausforderungen, die generische Webanwendungen in dieser Form nicht haben. Der Grund dafür ist, wie Karten und Geodaten im Browser dargestellt und geladen werden. Ein zentrales Problem ist die Darstellung der Karte selbst. OpenLayers rendert die Karte seit Version~3.18.0 standardmäßig über den Canvas-Renderer \parencite[v3.18.0]{openlayers_upgrade_2026} und zeichnet die Karteninhalte damit als Pixelgrafik statt als einzelne \ac{DOM}-Elemente. Ein Lokator kann zwar das Canvas-Element finden, aber nicht die Inhalte darin, beispielsweise ob ein bestimmter Layer gerade sichtbar ist. Was in der Karte dargestellt wird, ist also nicht direkt über das \ac{DOM} zugänglich. Dieses Problem ist nicht auf WebGIS-Anwendungen beschränkt, sondern gilt allgemein für Canvas-Elemente. Sie besitzen keinen repräsentativen \ac{DOM}-Baum für ihre innere Struktur, weshalb klassische, \ac{DOM}-basierte Verfahren darauf nicht anwendbar sind \parencite[S.~1]{bajammal_web_2018}.

Hinzu kommt, dass Geodaten asynchron über das Netz geladen werden. Kartenkacheln, \ac{WMS}-Layer, GetFeatureInfo-Abfragen oder Anfragen des Geocoders treffen mit unterschiedlicher Verzögerung ein. Ein Test muss abwarten, bis die Daten geladen und gerendert wurden, bevor getestet werden kann. Das in Abschnitt 2.2.2 beschriebene automatische Warten von Playwright hilft hier nur bedingt, da es sich auf die Bedienbarkeit von \ac{DOM}-Elementen bezieht. Ob Kartenkacheln vollständig dargestellt werden, lässt sich darüber nicht erkennen, sodass auf andere Methoden wie das Abwarten einer Netzwerkantwort zurückgegriffen werden muss.

Die Oberfläche einer WebGIS-Anwendung ist außerdem stark vom aktuellen Zustand abhängig. Ein Layer kann sichtbar oder unsichtbar sein, die Legende bezieht sich auf sichtbare Layer und eine Informationsanzeige erscheint erst, nachdem ein Feature in der Karte angeklickt wurde. Auch hier ist das Verhalten von Interaktionen wichtig. Die Informationsanzeige bei Klick auf ein Feature kann zum Beispiel deaktiviert sein, wenn gerade eine Messfunktion zum Messen einer Linie aktiv ist. Dieselbe Interaktion kann je nach Zustand zu einem anderen Ergebnis führen, weshalb die Vorbedingungen eines Tests genau definiert sein müssen.

Außerdem sind viele Interaktionen abhängig von Koordinaten. Angenommen, der Nutzer klickt auf eine Position, auf eine Markierung oder zeichnet durch mehrere Klicks eine Linie. Der Test muss dafür auf Pixelkoordinaten im Canvas klicken statt auf ein bekanntes Element. Das Ergebnis hängt dann von Kartenausschnitt, Zoomstufe und Projektion ab.